% ch07_u10 —— 原书 p.326–p.335（7.12.18 命题 -- 7.12.24 例 (d)）
% 底本：A. Beilinson, V. Drinfeld, Quantization of Hitchin's Integrable System
% and Hecke Eigensheaves，第 7 章。
% 承接：上单元（p.316–p.325）末句断在 "...are surjective"，本单元正文自「是满射」续起。
% 尾部接 p.336。

\SourcePage{326}

是满射\footnote{这比 $u_{\alpha\beta}$ 的满射性更强；例如，若 $M$ 是平坦 Mittag-Leffler $A$-模而非严格 Mittag-Leffler 模，则 7.12.6 的论证表明，$F_M$ 沿零截面的完备化不能表示为 $\varinjlim\operatorname{Spec}R_{\alpha}$，使映射 (356) 为满射。}，且理想 $\operatorname{Ker}u_{\alpha\beta}$ 幂零。
我们意义下的合理 $\aleph_0$-形式概形，就是 EGA I 意义下的形式概形。Grothendieck 意义下的拟紧形式概形，若具有由「defining ideals (English?)」（「Idéaux de définition」；见 EGA I 10.5.1）构成的基本系统，则是我们意义下的形式概形；特别地，EGA I 意义下的 noetherian 形式概形都如此。
% [Tutor Note: 原书 p.326 正文即印作「defining ideals (English?)」，括号中的 (English?) 是作者对英文术语的夹注，并非排版残留；底本如此，保留原样。]

由于我们主要关心域上 ind-finite type 的仿射 $\aleph_0$-形式概形，我们的定义与 EGA I 的定义之间的差别无关紧要。

\bdp{7.12.18} \emph{命题.} 设 $X$ 是 $A$ 上 ind-finite type 的形式光滑 $\aleph_0$-归纳概形，$S\subset X$ 是闭子概形，且 $S\to\operatorname{Spec}A$ 是同构。设 $X_{red}=S_{red}$（特别地，$X$ 是形式概形）。以 $M$ 记相对切层 $\Theta_{X/A}$ 限制到 $S$ 后的整体截面所成的 $A$-模。则 $M$ 是可数生成的投射模，且 $(X,S)$ 同构于归纳概形 $F_M$（见 7.12.14）沿零截面的完备化 $\widehat F_M$。

\emph{注.} 在 $A$ 上 ind-finite type 的形式光滑归纳概形 $X$ 上，$\mathcal{O}^{p}$-模 $\Theta_{X/A}$ 的定义与 $A=\mathbb{C}$ 的情形完全相同（见 7.11.8、7.11.7）。

\emph{证明.} 与 7.12.12 中一样，可证 $M$ 是平坦 Mittag-Leffler 模。$\aleph_0$ 假设蕴涵 $M$ 可数生成。由 7.12.8，$M$ 投射。

把 $X$ 表示为 $\varinjlim X_n$（$n\in\mathbb{N}$），其中 $X_n$ 是 $X$ 的含 $S$ 的闭子概形，且 $X_n\subset X_{n+1}$。设 $X^{(1)}$ 是 $S$ 在 $X$ 中的第一次无穷小邻域，即 $X^{(1)}$ 是 $S$ 在 $X_n$ 中的第一次无穷小邻域之并，
\SourcePage{327}%
$n\in\mathbb{N}$。显然 $X^{(1)}=F^{(1)}_{M}:=0\in F_M$ 的第一次无穷小邻域。

嵌入 $X^{(1)}\to F_M$ 可延拓为态射 $\varphi:X\to\widehat F_M$（为构造 $\varphi$，先构造 $\varphi_n:X_n\to\widehat F_M$，使 $\varphi_n|_{X_{n-1}}=\varphi_{n-1}$，且 $\varphi_n$ 在 $X_n\cap X^{(1)}$ 上的限制是典范嵌入 $X_n\cap X^{(1)}\hookrightarrow F^{(1)}_{M}$；这之所以可行，是因为 $\widehat F_M$ 在 $A$ 上形式光滑）。完全类似地，把嵌入 $F^{(1)}_{M}=X^{(1)}\hookrightarrow X$ 延拓为态射 $\psi:\widehat F_M\to X$。因 $\varphi$ 与 $\psi$ 诱导 $F^{(1)}_{M}$ 与 $X^{(1)}$ 之间的同构，故 $\varphi$ 与 $\psi$ 都是 ind-闭嵌入，且 $\varphi\psi$ 是同构。于是 $\varphi$ 与 $\psi$ 是同构。

\bdp{7.12.19} \emph{例.} 我们将构造一对 $(X,S)$，它满足 7.12.18 中除 $\aleph_0$ 假设以外的全部条件，且 $(X,S)$ 不 $A$-同构于形如 $\widehat F_M$ 的形式概形。

设给定平坦 Mittag-Leffler 模的一个非平凡扩张
\begin{equation}\tag{357}
0\to N'\to N\to L\to 0.
\end{equation}
这样的扩张对「大多数」环 $A$ 都存在；见 7.12.24(b, a$''$, d)。把 (357) 与 $A[t]$ 张量，得到扩张 $0\to N'[t]\to N[t]\to L[t]\to 0$。把此扩张乘以 $t$，得到 $0\to N'[t]\to Q\to L[t]\to 0$。归纳概形 $F_Q$ 在 $A[t]$ 上形式光滑，从而在 $A$ 上形式光滑。设 $S\subset F_Q$ 是零截面复合 $\operatorname{Spec}A\to\operatorname{Spec}A[t]\to F_Q$ 的像。以 $X$ 记 $F_Q$ 沿 $S$ 的完备化。

在证明 $(X,S)$ 具有所需性质之前，我们先更具体地描述 $X$。对 $A$-代数 $R$，$F_Q$ 的一个 $R$-点由两部分构成：一个 $A$-态射 $A[t]\to R$，以及 $Q\otimes_{A[t]}R$ 的一个元素。换言之，$F_Q$ 的 $R$-点由一个三元组 $(n,l,t)$ 给出，其中 $n\in N\otimes_A R$，$l\in L\otimes_A R$，$t\in R$，且
\begin{equation}\tag{358}
\pi(n)=tl,
\end{equation}
这里 $\pi$ 是投影 $N\otimes_A R\to L\otimes_A R$。

\SourcePage{328}

于是 $F_Q$ 是 $F_N\times F_L\times\mathbb{A}^1$ 的闭归纳子概形，由方程 (358) 定义。因此 $X\subset\widehat F_N\times\widehat F_L\times\widehat{\mathbb{A}}^1$ 由同一方程 (358) 定义（这里 $\widehat{\mathbb{A}}^1$ 是 $\mathbb{A}^1$ 在 $0\in\mathbb{A}^1$ 处的完备化）。

现设 $(X,S)$ 与 $\widehat F_M$ 是 $A$-同构的。则 $M$ 就是 $\Theta_{X/A}$ 限制到 $S$ 后的整体截面所成的模。把 (358) 线性化，得到
\begin{equation}\tag{359}
M=N'\oplus L\oplus A\subset N\oplus L\oplus A.
\end{equation}
复合
\begin{equation}\tag{360}
\widehat F_M\xrightarrow{\sim}X\hookrightarrow\widehat F_N\times\widehat F_L\times\mathbb{A}^1
\end{equation}
由「Taylor 级数」$\sum_{n=1}^{\infty}\varphi_n$ 给出，其中 $\varphi_n$ 是次数为 $n$ 的齐次多项式映射 $M\to N\oplus L\oplus A$；显然 $\varphi_1$ 是嵌入 (359)。置 $f=\operatorname{pr}_N\circ\varphi_2$，其中 $\operatorname{pr}_N$ 是投影 $N\oplus L\oplus A\to N$。因 $M=N'\oplus L\oplus A$，二次映射 $M\to N$ 构成的模以双线性映射 $L\times A\to N$（即 $\operatorname{Hom}(L,N)$）构成的模为直和项。$f$ 在 $\operatorname{Hom}(L,N)$ 中的像给出 (357) 的一个分裂（用下述事实：态射 (360) 经过由方程 (358) 定义的归纳子概形 $X\subset\widehat F_N\times\widehat F_L\times\mathbb{A}^1$ 分解）。于是得到矛盾。

\bdp{7.12.20} \emph{命题.} 设 $X$ 是环 $A$ 上形式光滑的归纳概形。设下列两个假设之一成立：(i) $X$ 是 ind-affine 的；(ii) $A$ noetherian，且 $X$ 在 $A$ 上 ind-finite type。则 $X$ 是一族在 $A$ 上形式光滑的 ind-闭 $\aleph_0$-归纳概形的有向并。

\emph{证明.} 只需证明：对闭子概形的任意递增序列 $Y_n\subset X$，存在在 $A$ 上形式光滑的 ind-闭 $\aleph_0$-归纳概形 $Y\subset X$，使 $Y\supset Y_n$ 对所有 $n$ 成立。

设 $X$ 是 ind-affine 的。则每个 $Y_n$ 都是仿射的。把 $Y_n$ 表示为 $A$ 上形式光滑概形 $V_n$ 的闭子概形（例如，把 $Y_n$ 的坐标环表示为 $A$ 上多项式代数的商）。

\SourcePage{329}

设 $Y'_n\subset V_n$ 是 $Y_n$ 在 $V_n$ 中的第一次无穷小邻域。因 $X$ 形式光滑，态射 $Y_n\hookrightarrow X$ 可延拓为态射 $Y'_n\to Z_n\subset X$，其中 $Z_n\subset X$ 是某个闭子概形。置 $Y^{(2)}_n:=Z_1\cup\ldots\cup Z_n$。现把上述构造用于 $(Y^{(2)}_n)$，得到新序列 $(Y^{(3)}_n)$，如此类推。所有 $Y^{(k)}_n$ 之并在 $A$ 上形式光滑。

若 $X$ ind-quasicompact 但非 ind-affine，则上述构造的显然修改给出含所有 $Y_n$ 的 ind-闭 $\aleph_0$-归纳概形 $Y\subset X$，使得对 $A$ 上任意仿射概形 $S$ 与由理想 $\mathcal{I}\subset\mathcal{O}_S$（$\mathcal{I}^2=0$）定义的闭子概形 $S_0\subset S$，每个 $A$-态射 $S_0\to Y$ 都局部地延拓为态射 $S\to Y$。若假设 (ii) 成立，则这蕴涵整体延拓的存在。

\bdp{7.12.21} 我们将描述域 $C$ 上 ind-finite type 的形式光滑仿射 $\aleph_0$-形式概形（据 7.12.20，一般情形在某种意义上可归到 $\aleph_0$ 情形）。首先有如下例子。

(0) 置 $R_{mn}:=C[x_1,\ldots,x_m][[x_{m+n},\ldots,x_{m+n}]]$。则 $\operatorname{Spf}R_{mn}$ 是 $C$ 上形式光滑的仿射 $\aleph_0$-形式概形。
% [Tutor Note: 原书 p.329 此处印作 $C[x_1,\ldots,x_m][[x_{m+n},\ldots,x_{m+n}]]$，上下标重复「m+n」；据上下文及原书 p.330 的同一记号，应为 $[[x_{m+1},\ldots,x_{m+n}]]$。底本如此，保留原样。]

(i) 设 $I\subset R_{mn}$ 是理想，$A:=R_{mn}/I$。以 $\mathcal{I}$ 记 $\operatorname{Spf}R_{mn}$ 上对应于 $I$ 的理想层。当然，$\operatorname{Spf}A$ 是 $C$ 上 ind-finite type 的仿射 $\aleph_0$-形式概形。它形式光滑，当且仅当对每个 $u\in\operatorname{Spf}A$，$\mathcal{I}$ 在 $u$ 处的茎由某些 $f_1,\ldots,f_r\in I$ 生成，且 Jacobi 矩阵 $\bigl(\frac{\partial f_i}{\partial x_j}(u)\bigr)$ 的秩为 $r$。

(ii) 设 $A$ 如 (i)，且 $\operatorname{Spf}A$ 形式光滑。则 $\operatorname{Spf}A[[y_1,y_2,\ldots]]$ 是 $C$ 上 ind-finite type 的形式光滑仿射 $\aleph_0$-形式概形。

在 7.12.22 与 7.12.23 中我们将证明：域上每个连通的、ind-finite type 的形式光滑仿射 $\aleph_0$-形式概形，都同构于例 (i) 或 (ii) 中的形式概形。

\SourcePage{330}

\bdp{7.12.22} \emph{命题.} 设 $X$ 是域 $C$ 上 ind-finite type 的形式光滑的仿射形式概形，且 $\Theta_X$ 凝聚（即 $\Theta_X$ 到 $X$ 的每个闭子概形的限制都有限生成）。则 $X$ 同构于例 7.12.21(i) 中的形式概形。

\emph{证明.} 把 $X$ 表示为 $\varinjlim\operatorname{Spec}A_i$，使当 $i\le j$ 时态射 $A_j\to A_i$ 是核幂零的满射。各代数 $A_i$ 都是有限型的。可设指标集 $i$ 有最小元 $0$。置 $I_i:=\operatorname{Ker}(A_i\to A_0)$。

\emph{引理.} 对每个 $k\in\mathbb{N}$，存在 $i_1$，使态射 $A_i/I_i^k\to A_{i_1}/I_{i_1}^k$ 对一切 $i\ge i_1$ 都是双射。

假定引理，对足够大的 $i$ 置 $A_{(k)}:=A_i/I_i^k$，$I_{(k)}:=\operatorname{Ker}(A_{(k)}\to A_0)$。显然 $A_{(1)}=A_0$，$A_{(k)}=A_{(k+1)}/I^k_{(k+1)}$，$I_{(k)}=I_{(k+1)}/I^k_{(k+1)}$。有 $X=\operatorname{Spf}A$，其中 $A:=\varprojlim A_{(k)}$。取 $A_{(1)}=A_0$ 的生成元 $\bar x_1,\ldots,\bar x_m$，以及 $A_0$-模 $I_{(2)}$ 的生成元 $\bar x_{m+1},\ldots,\bar x_{m+n}$。把 $\bar x_1,\ldots,\bar x_{m+n}$ 提升为 $\tilde x_1,\ldots,\tilde x_{m+n}\in A$。置
\[
R_{mn}:=C[x_1,\ldots,x_m][[x_{m+1},\ldots,x_{m+n}]].
\]
存在唯一的连续同态 $f:R_{mn}\to A$，使 $x_i\mapsto\tilde x_i$。显然 $f$ 是满射。此外，$f$ 诱导满射 $\mathfrak{a}^k\twoheadrightarrow\operatorname{Ker}(A\to A_{(k)})$，其中 $\mathfrak{a}\subset R_{mn}$ 是由 $x_{m+1},\ldots,x_{m+n}$ 生成的理想。故 $f$ 是开映射。因此 $f$ 诱导 $A$ 与 $R_{mn}$ 的某个商之间的拓扑同构。命题得证。

余下证明引理。存在 $i_0$，使对每个 $i\ge i_0$，态射 $\operatorname{Spec}A_{i_0}\to\operatorname{Spec}A_i$ 诱导切空间之间的同构（事实上，因 $\Theta_X$ 到 $\operatorname{Spec}A_0$ 的限制有限生成，对应于 $A_0$-模 $N_i:=\Omega_i\otimes_{A_i}A_0$ 的函子 (355) 同构于函子 $L\mapsto\operatorname{Hom}(Q,L)$，其中 $Q$ 是某个 $A_0$-模；故存在 $i_0$，使 $N_i=N_{i_0}$ 对 $i\ge i_0$ 成立）。可设 $i_0=0$。置 $Y_i:=\operatorname{Spec}A_i/I_i^k$（特别地，$Y_0=\operatorname{Spec}A_0$）。态射 $Y_0\to Y_i$ 诱导切空间之间的同构。

\SourcePage{331}

把 $A_0$ 表示为 $C[x_1,\ldots,x_n]/J$，并置 $\tilde Y_0:=\operatorname{Spec}C[x_1,\ldots,x_n]/J^k$。因 $X$ 形式光滑，态射 $Y_0\hookrightarrow X$ 可延拓为态射 $\tilde Y_0\to X$。其像含于某个 $Y_{i_1}$ 中。我们来证明：对 $i\ge i_1$，嵌入 $\nu:Y_{i_1}\hookrightarrow Y_i$ 是同构。我们有态射 $f:\tilde Y_0\to Y_{i_1}$。另一方面，态射 $Y_0\hookrightarrow\tilde Y_0$ 延拓为 $g:Y_i\to\tilde Y_0$。复合 $\nu fg:Y_i\to Y_i$ 在 $Y_0$ 上诱导恒等。故 $\nu fg$ 有限，且诱导切空间之间的同构。因此 $\nu fg$ 是闭嵌入。因 $Y_i$ noetherian，闭嵌入 $Y_i\to Y_i$ 是同构。故 $\nu fg$ 是同构，从而 $\nu$ 是同构。

\bdp{7.12.23} \emph{命题.} 设 $X$ 是域 $C$ 上连通的、ind-finite type 的形式光滑仿射 $\aleph_0$-形式概形，且 $\Theta_X$ 不凝聚（即 $\Theta_X$ 到 $X_{red}$ 的限制是无限型的）。则 $X$ 同构于例 7.12.21(ii) 中的形式概形。

\emph{证明.} 我们将构造形式光滑态射
\[
X\to\operatorname{Spf}C[[y_1,y_2,\ldots]],
\]
它在 $0\in\operatorname{Spf}C[[y_1,y_2,\ldots]]$ 上的纤维是 7.12.21(i) 中的形式概形。

把 $X$ 表示为 $\varinjlim\operatorname{Spec}A_n$（$n\in\mathbb{N}$），使对每个 $n$，态射 $A_{n+1}\to A_n$ 是核幂零的满射。各代数 $A_n$ 都是有限型的。由 7.12.13，$\Theta_X$ 到 $\operatorname{Spec}A_n$ 的限制自由；其秩可数。这意味着：对每个 $n$，投射系统 $(\Omega_{A_i}\otimes_{A_i}A_n)$（$i\ge n$）等价于投射系统
\[
\ldots\to A_n^3\to A_n^2\to A_n
\]
（这里映射 $A_n^{k+1}\to A_n^k$ 是到前 $k$ 个坐标的投影）。故把序列 $(A_n)$ 换成其子序列后，得到图解
\[
\ldots\twoheadrightarrow\Omega_{A_3}\twoheadrightarrow F_2\twoheadrightarrow\Omega_{A_2}\twoheadrightarrow F_1\twoheadrightarrow\Omega_{A_1}
\]
\SourcePage{332}
其中 $F_n$ 是有限生成的自由 $A_n$-模，且 $A_n$-模 $G_n:=\operatorname{Ker}(F_{n+1}\otimes_{A_{n+1}}A_n\to F_n)$ 也自由。对每个 $n$ 取 $G_n$ 的一组基 $e_{n1},\ldots,e_{nk_n}$。把 $e_{ni}$ 提升为 $\tilde e_{ni}\in\operatorname{Ker}(\Omega_{A_{n+2}}\otimes_{A_{n+2}}A_n\to F_n)\subset\operatorname{Ker}(\Omega_{A_{n+2}}\otimes_{A_{n+2}}A_n\to\Omega_{A_n})$，并把 $\tilde e_{ni}$ 表示为 $df_{ni}$，其中 $f_{ni}\in\operatorname{Ker}(A_{n+2}\to A_2)$。
% [Tutor Note: 原书 p.332 此处印作 $f_{ni}\in\operatorname{Ker}(A_{n+2}\to A_2)$，按上下文应为 $A_n$（$\tilde e_{ni}$ 是 $\Omega_{A_{n+2}}\otimes_{A_{n+2}}A_n$ 中的元素，$df_{ni}$ 的 $f_{ni}$ 应在 $A_{n+2}$ 中而像在 $A_n$ 中）。底本如此，保留原样。]
最后把 $f_{ni}$ 提升为 $\tilde f_{ni}\in A:=\varprojlim_m A_m$，并把 $f_{ni}$（$n\in\mathbb{N}$，$i\le k_n$）排成一个序列 $\varphi_1,\varphi_2,\ldots$。此序列收敛到 $0$，故有连续同态 $C[[y_1,y_2,\ldots]]\to A$，使 $y_i\mapsto\varphi_i$。它诱导态射
\begin{equation}\tag{361}
f:X\to Y:=\operatorname{Spf}C[[y_1,y_2,\ldots]].
\end{equation}
由构造可见，微分
\begin{equation}\tag{362}
df:\Theta_X\to f^*\Theta_Y
\end{equation}
是满射，且其核凝聚（事实上，(362) 到 $\operatorname{Spec}A_1\subset X$ 的限制显然具有这些性质，故到 $\operatorname{Spec}A_n$（$n\in\mathbb{N}$）的限制也具有这些性质）。

\emph{引理.} 形式光滑、ind-finite type 的归纳概形之间的态射 $f:X\to Y$ 形式光滑，当且仅当其微分 (362) 是满射。此时 $\Theta_{X/Y}$ 是 (362) 的核。

假定引理，可见 (361) 形式光滑，且 $\Theta_{X/Y}$ 凝聚。故 (361) 在 $0\in Y$ 上的纤维 $X_0$ 满足命题 7.12.22 的条件。因此 $X_0$ 同构于例 7.12.21(i) 中的形式概形。我们来证明 $X$ 同构于 $\tilde X:=X_0\times Y$。因 $X$ 在 $Y$ 上形式光滑，嵌入 $X_0\hookrightarrow X$ 延拓为 $Y$-态射 $\alpha:\tilde X\to X$。因 $\tilde X$ 在 $Y$ 上形式光滑，嵌入 $X_0\hookrightarrow\tilde X$ 延拓为 $Y$-态射 $\beta:X\to\tilde X$。$\alpha$ 与 $\beta$ 都是 ind-闭嵌入（若有限型概形的态射 $\nu:Y\to Z$ 诱导同构 $Y_{red}\to Z_{red}$，且 $\nu$ 的每个几何纤维既约，则 $\nu$ 是闭嵌入）。$Y$-态射 $\beta\alpha:X_0\times Y\to X_0\times Y$
\SourcePage{333}%
在 $0\in Y$ 上诱导恒等，故 $\beta\alpha$ 是同构。因此 $\alpha$ 与 $\beta$ 是同构，命题得证。

引理的证明是标准的。关于 $\Theta_{X/Y}$ 的陈述由定义推出。为证第一个陈述，取仿射概形 $S$ 及理想 $\mathcal{I}\subset\mathcal{O}_S$（$\mathcal{I}^2=0$），以 $S_0\subset S$ 记对应于 $\mathcal{I}$ 的子概形。对态射 $\psi:S_0\to X$，以 $E_X(S,\mathcal{I},\psi)$（相应地 $E_Y(S,\mathcal{I},\psi)$）记 $\psi$（相应地 $f\psi$）到态射 $S\to X$（相应地 $S\to Y$）的延拓之集。$f$ 的形式光滑性指：对上述所有 $S,\mathcal{I},\psi$，$f_*:E_X(S,\mathcal{I},\psi)\to E_Y(S,\mathcal{I},\psi)$ 是满射。因 $X$ 与 $Y$ 形式光滑，$E_X(S,\mathcal{I},\psi)$ 与 $E_Y(S,\mathcal{I},\psi)$ 非空。据 [Gr67] 的 16.5.14，它们分别是 $V_X(S,\mathcal{I},\psi)$ 与 $V_Y(S,\mathcal{I},\psi)$ 上的 torsor（即非空的仿射空间），其中 $V_X(S,\mathcal{I},\psi):=\operatorname{Hom}(\psi^*\Omega_X,\mathcal{I})=\Gamma(S_0,\psi^*\Theta_X\otimes\mathcal{I})$，$V_Y(S,\mathcal{I},\psi)=\Gamma(S_0,\psi^*f^*\Theta_Y\otimes\mathcal{I})$。映射 $f_*$ 是仿射的，而相应的线性映射 $\Gamma(S_0,\psi^*\Theta_X\otimes\mathcal{I})\to\Gamma(S_0,\psi^*f^*\Theta_Y\otimes\mathcal{I})$ 由 (362) 诱导。故引理的第一个陈述显然成立。

\bdp{7.12.24} \emph{Mittag-Leffler 模的例子.}

(a) 据 [RG] p.77 的 2.4.1，对每个 noetherian $A$ 与投射 $A$-模 $P$，$A$-模 $P^*:=\operatorname{Hom}_A(P,A)$ 是严格 Mittag-Leffler 且平坦的。为证 $P^*$ 严格 Mittag-Leffler，可如下论证：对任意 $f:F\to P^*$（$F$ 有限型），$f^*:P\to F^*$ 的像由某些 $l_1,\ldots,l_n\in F^*$ 生成；这些 $l_i$ 给出 $u:F\to A^n$，使 $f=gu$ 且 $u=hf$，其中 $g:A^n\to P^*$，$h:P^*\to A^n$。

特别地，若 $A$ noetherian，则对每个集合 $I$，$A$-模 $A^I$ 是严格 Mittag-Leffler 且平坦的。

(a$'$) 众所周知，若 $A$ 是 Dedekind 环而非域，则对无限集 $I$，$A^I$ 不投射。事实上，可设 $I$ 可数。固定非零素理想 $\mathfrak{p}\subset A$，考虑由满足 $a_i\to 0$（在 $\mathfrak{p}$-进拓扑下）的元素 $a=(a_i)\in A^I$ 构成的子模 $M$。若 $A^I$ 投射，则局部化 $M_{\mathfrak{p}}$ 自由。因 $M/\mathfrak{p}M$ 的维数可数，$M_{\mathfrak{p}}$ 的秩应可数。但 $M$ 含一个同构于 $A^I$ 的子模，故 $(A^I)_{\mathfrak{p}}$ 的秩应可数。然而 $(A^I)_{\mathfrak{p}}/\mathfrak{p}\cdot(A^I)_{\mathfrak{p}}=(A/\mathfrak{p})^I$ 的维数不可数，矛盾。

\SourcePage{334}

(a$''$) 设 $A$ 在 $\mathbb{Z}$ 上或在域上有限生成\footnote{我们不知道是否只需假定 $A$ noetherian。}。若 $A$ 不是 Artin 环，且 $I$ 无限，则 $A^I$ 不投射：用 (a$'$) 以及 $A$ 上有限 Dedekind 环 $B$ 的存在性。

(b) 若 $L$ 是非投射的平坦 Mittag-Leffler 模，则存在非分裂正合列 $0\to N'\to N\to L\to 0$，其中 $N$ 与 $N'$ 是平坦 Mittag-Leffler 模。事实上，若 $N$ 是投射模，$N\to L$ 是满射，则它不分裂，且 $\operatorname{Ker}(N\to L)$ 是 Mittag-Leffler 的（[RG] p.71 的 2.1.6）。

(c) [RG] 中指出：若
\begin{equation}\tag{363}
0\to A\xrightarrow{f}M'\to M\to 0
\end{equation}
是 $A$-模的非分裂正合列，且 $M$ 平坦、Mittag-Leffler，则 $M'$ 是 Mittag-Leffler 的，但不是严格 Mittag-Leffler 的。事实上，若 $M'$ 严格 Mittag-Leffler，则存在有限表现模 $G$ 与态射 $u:A\to G$，使 $f=gu$ 且 $u=hf$，其中 $g:G\to M'$，$h:M'\to G$。因 $M$ 是有限生成投射模的直极限，可设 $\operatorname{Im}g\subset\operatorname{Im}f$。于是 $gh$ 给出 (363) 的分裂，即得矛盾。

这里还有另一个论证。$F_{M'}$ 在 $0\in F_M$ 上的纤维是 $F_{M'}$ 的闭子概形，典范同构于 $\operatorname{Spec}A\times\mathbb{A}^1$；若 (363) 非分裂，则投影 $\operatorname{Spec}A\times\mathbb{A}^1\to\mathbb{A}^1$ 不能延拓为函数 $F_{M'}\to\mathbb{A}^1$，故由 7.12.16，$M'$ 不是严格 Mittag-Leffler 的。

\SourcePage{335}

(d) 设 $A$ 是 Dedekind 环，既非域也非完备局部环。则据 [RG] p.76，存在非分裂正合列 (363)，其中 $M$ 是平坦的严格 Mittag-Leffler $A$-模。下面是它的构造。以 $K$ 记 $A$ 的分式域。固定非零素理想 $\mathfrak{p}\subset A$，考虑完备化 $\widehat A_{\mathfrak{p}}$、$\widehat K_{\mathfrak{p}}$；则 $\widehat A_{\mathfrak{p}}\ne A$，$\widehat K_{\mathfrak{p}}\ne K$。以 $M$ 记满足 $a_n\in\mathfrak{p}^{-n}$ 且在 $\widehat K_{\mathfrak{p}}$ 中收敛的序列 $(a_n)$ 构成的模；有态射 $\lim:M\to\widehat K_{\mathfrak{p}}$。注意 $M$ 是严格 Mittag-Leffler 模\footnote{$M$ 是 Mittag-Leffler 模这一事实是显然的：$A$ 是 Dedekind 环，$M$ 平坦，且对每个有限维子空间 $V\subset M\otimes K$，模 $V\cap M$ 有限生成}。事实上，据上面的 (a)，$\prod_{n=1}^{\infty}\mathfrak{p}^{-n}$ 是严格 Mittag-Leffler 的，而 $\bigl(\prod_{n=1}^{\infty}\mathfrak{p}^{-n}\bigr)/M$ 平坦，故 $M$ 严格 Mittag-Leffler。我们断言 $\operatorname{Ext}(M,A)\ne 0$，即态射 $\varphi:\operatorname{Hom}(M,K)\to\operatorname{Hom}(M,K/A)$ 不满。更确切地说，设 $l:M\to K/A$ 是 $\lim:M\to\widehat K_{\mathfrak{p}}$ 与态射 $\widehat K_{\mathfrak{p}}\to\widehat K_{\mathfrak{p}}/\widehat A_{\mathfrak{p}}\hookrightarrow K/A$ 的复合。我们将证明 $l\notin\operatorname{Im}\varphi$。

设 $l$ 来自 $\tilde l:M\to K$。$\tilde l$ 在 $\mathfrak{p}^{-n}\subset M$ 上的限制定义 $c_n\in\operatorname{Hom}(\mathfrak{p}^{-n},A)=\mathfrak{p}^n$。则 $\tilde l=\tilde l'$，其中 $\tilde l':M\to K_{\mathfrak{p}}$ 把 $(a_n)\in M$ 送到
% [Tutor Note: 原书 p.335 此处印作 $\tilde l':M\to K_{\mathfrak{p}}$（$K$ 上无 hat），而前文 $\lim:M\to\widehat K_{\mathfrak{p}}$；按上下文应为 $\widehat K_{\mathfrak{p}}$。底本如此，保留原样。]
\begin{equation}\tag{364}
\sum_{n=1}^{\infty}c_na_n+\lim_{n\to\infty}a_n.
\end{equation}
事实上，$\tilde l'-\tilde l$ 是态射 $M/M_0\to\widehat A_{\mathfrak{p}}$，其中 $M_0$ 是使 $a_n=0$ 对足够大的 $n$ 成立的 $(a_n)\in M$ 之集；另一方面，$\operatorname{Hom}(M/M_0,\widehat A_{\mathfrak{p}})=0$，因为 $M/M_0$ 是 $\mathfrak{p}$-可除的（即 $\mathfrak{p}M+M_0=M$）。因 $\tilde l'=\tilde l$，表达式 (364) 对每个序列 $(a_n)\in M$ 都属于 $K\subset\widehat K_{\mathfrak{p}}$。这是不可能的（分别考虑非零 $c_n$ 的个数有限与无限两种情形）。

% 尾部接 p.336
